\documentclass[11pt]{article}

\usepackage[a4paper,margin=1in]{geometry}
\usepackage{amsmath}
\usepackage{amssymb}
\usepackage{graphicx}
\usepackage{booktabs}
\usepackage{array}
\usepackage{caption}
\usepackage{hyperref}

\hypersetup{
  colorlinks=true,
  linkcolor=black,
  urlcolor=blue!60!black,
  pdftitle={INDENG 231 Project 1},
  pdfauthor={Arthur Fang}
}

\setlength{\parindent}{0pt}
\setlength{\parskip}{0.55em}
\renewcommand{\arraystretch}{1.1}
\captionsetup{font=small, labelfont=bf, skip=4pt}
\graphicspath{{../outputs/figures/}}

\begin{document}

\noindent\textbf{Arthur Fang} \hfill INDENG 231 Project 1 \hfill April 29, 2026

\vspace{0.6em}\hrule\vspace{0.6em}

% ===============================================================
\section{Backtesting System Design}

\textbf{Data layer.} \texttt{src/data\_loader.py} loads close-price data
into a date $\times$ ticker DataFrame; both wide and long CSV layouts are
auto-detected. Missing prices are forward-filled and tickers that remain
entirely NaN are dropped. Daily simple returns
$R_{i,t} = P_{i,t}/P_{i,t-1} - 1$ live in a parallel panel of identical
shape.

\textbf{Strategy contract.} \texttt{src/strategies/base.py} defines
\texttt{BaseStrategy.generate\_weights(prices\_until\_t, returns\_until\_t,
current\_date)}, returning a target weight Series for the current day.
Subclasses override that single method; the engine handles everything
else. A shared \texttt{validate\_weights} method clips negatives to zero,
normalizes if the gross exceeds 1, and treats the residual as cash. Three
reusable weighting helpers
(\texttt{equal\_weight}, \texttt{inverse\_volatility\_weight},
\texttt{rank\_weight}) live in \texttt{src/portfolio.py}.

\textbf{Engine.} \texttt{src/backtester.py} iterates the trading calendar.
At day $t$ it slices prices/returns to rows \texttt{[: t]}, asks the
strategy for $w_t$, validates, and \emph{realizes those weights against
$R_{i,t+1}$:}
\[
   R_{p,\,t+1} = \sum_i w_{i,t} R_{i,t+1} - \kappa\,\lVert w_t - w_{t-1}\rVert_1,
   \qquad
   \mathrm{NAV}_{t+1} = \mathrm{NAV}_t (1 + R_{p,\,t+1}).
\]
The strategy never sees $R_{i,t+1}$ when choosing $w_{i,t}$; this is the
structural no-lookahead guarantee. An optional
\texttt{update\_after\_return} hook lets adaptive strategies (the UCB
meta-strategy in \S\ref{sec:ext}) receive realized rewards, and is fired
\emph{after} $R_{p,t+1}$ is computed.

\textbf{Metrics.} \texttt{src/metrics.py} computes cumulative return,
annualized return, annualized volatility, Sharpe (zero risk-free rate),
maximum drawdown, Calmar, win rate, and average L1 turnover. All
annualizations use $T = 252$.

\textbf{Outputs.} Each phase writes one CSV (\texttt{outputs/tables/}),
two PNGs (\texttt{outputs/figures/}, NAV + drawdown overlays), and one
JSON run log (\texttt{outputs/logs/}).

\textbf{Assumptions.} (i) Weights execute at the recorded close;
(ii) zero risk-free rate; (iii) long-only, no leverage; (iv) transaction
cost defaults to 0 (the engine accepts a proportional $\kappa$ for
sensitivity); (v) no corporate-action handling, prices treated as
adjusted; (vi) static universe equal to the loaded CSV (survivorship-biased).

\textbf{Module map.} \texttt{src/} (config, data\_loader, backtester,
metrics, plotting, portfolio, risk), \texttt{src/strategies/} (base,
single\_stock, benchmarks, cross\_sectional, bandit),
\texttt{experiments/} (one entry script per deliverable),
\texttt{outputs/}, and \texttt{report/}.

% ===============================================================
\section{Design Choices}

\textbf{Next-day return realization, not shifted weights.}
The strategy sees \texttt{prices[:t]} (close of $t$ inclusive) and is
paid by $R_{i,t+1}$. The alternative --- pay $w_t$ with $R_{i,t}$ ---
would force the strategy to choose with information through $t-1$,
which conflates ``observe the close'' with ``trade on the close''. The
chosen pattern matches how a real desk works: decide at the close,
realize the move overnight to the next close.

\textbf{Constraint enforcement in the base class.}
\texttt{validate\_weights} clips negatives, normalizes if gross $>1$,
and asserts the invariants. Pushing this into the engine means a new
strategy author cannot accidentally short or lever, and the engine
never has to trust the subclass.

\textbf{Positive-return gate on the new momentum strategies.}
Restricting selection to names with $\rho_{i,t} > 0$ converts top-$K$
momentum into a long-or-cash rule for free: in a broad-market sell-off
the entire cross-section can be negative and the strategy parks in
cash instead of loading on the ``least bad'' names.

\textbf{$K{=}3$ with inverse-volatility weighting.}
A concentrated portfolio is dominated by idiosyncratic vol, so
inverse-vol weighting matters more there than at $K{=}10$ where the
spread is already diluted. The grid agreed with this intuition ---
the inverse-vol variant won at $K{=}3$ but only marginally at $K{=}10$.
Sitting at the boundary ($K{=}3$, $\ell{=}90$) is flagged in \S5 as
an overfitting risk; with this sample length the in-sample optimum
clearly wants \emph{more} concentration and \emph{longer} memory than
the grid offers, but $K{=}1$ is no longer a portfolio strategy.

\textbf{Drawdown threshold $-10\%$ on a 120-day window.}
$-5\%$ triggers on routine intra-month volatility (too many false
positives), $-15\%$ only fires after the worst is already past
(too late). 120 trading days $\approx 6$ months gives a smooth
medium-term drawdown signal that does not whipsaw on single-week
sell-offs. The grid confirmed $-10\%$ as the in-sample winner.

\textbf{60/40 split rather than walk-forward.}
A single split is transparent and easy to audit, which is the right
trade-off for a course project. Walk-forward retuning would reduce
overfitting more, at the cost of an extra hyperparameter (refit
frequency) and a much heavier compute bill --- worth doing as a
follow-up but not the headline protocol here.

% ===============================================================
\section{Deliverable 3 -- Single-Stock Strategy Comparison}

Test five long-or-cash strategies on NVDA over the full 1{,}255-day
sample: 20-day momentum; 20-day mean reversion (long after a $-5\%$
trailing return); 20/50 SMA crossover; volatility-filtered momentum
(long iff trailing return $> 0$ AND 20-day vol $< 3\%$); and 20-day
z-score mean reversion (long iff $z_t < -1$). Each strategy is a
single-stock specialization of the generic \texttt{BaseStrategy} ---
no engine code changes.

\begin{table}[h!]
\centering
\caption{Single-stock (NVDA) backtest metrics, full sample.}
\begin{tabular}{lrrrrrr}
\toprule
\textbf{Strategy} & \textbf{Cum.\ Ret.} & \textbf{Ann.\ Ret.} & \textbf{Ann.\ Vol.} & \textbf{Sharpe} & \textbf{Max DD} & \textbf{Turnover} \\
\midrule
Momentum (20)              & 3.694 & 0.364 & 0.384 & \textbf{0.997} & $-0.498$           & 0.095 \\
Z-score mean reversion (20)& 2.022 & 0.249 & 0.262 & 0.977          & $\mathbf{-0.226}$  & 0.110 \\
Vol-filtered momentum      & 2.131 & 0.258 & 0.272 & 0.974          & $-0.364$           & 0.076 \\
MA crossover (20/50)       & 3.042 & 0.324 & 0.376 & 0.930          & $-0.636$           & \textbf{0.022} \\
Mean reversion ($-5\%$)    & 1.838 & 0.233 & 0.298 & 0.851          & $-0.407$           & 0.075 \\
\bottomrule
\end{tabular}
\end{table}

\begin{figure}[h!]
\centering
\includegraphics[width=0.78\linewidth]{single_stock_nav.png}
\caption{Single-stock NAV curves on NVDA.}
\label{fig:nvda-nav}
\end{figure}

\begin{figure}[h!]
\centering
\includegraphics[width=0.78\linewidth]{single_stock_drawdown.png}
\caption{Single-stock drawdown profiles on NVDA.}
\label{fig:nvda-dd}
\end{figure}

\textbf{Interpretation.} The top three Sharpes (momentum 0.997,
z-score 0.977, vol-filtered 0.974) are within 0.02 of each other --- well
inside the noise band of a 5-year single-asset Sharpe estimate (the
approximate standard error of an annualized Sharpe at $N{=}1255$ is
$\approx 0.45$; see \S\ref{sec:ext}). Momentum captures the most return
($3.7\times$) but pays a $-49.8\%$ drawdown; z-score mean reversion
spends most of the period in cash and gives back roughly half the return
for less than half the drawdown, which is why it tops the Calmar ranking.
MA crossover is the lowest-turnover rule by an order of magnitude.

% ===============================================================
\section{Deliverable 4 -- Portfolio-Level Backtest}

Four constructions are tested over the full 101-stock pool. Two pin the
selection rule to the top-10 names by trailing 30-day return and vary
\emph{only} the weighting scheme, isolating the uniform vs.\ inverse-vol
comparison required by the rubric. The SMA crossover portfolio
(equal-weighted) and the equal-weight buy-and-hold of the entire pool
are included as anchors. Inverse-volatility uses a 20-day rolling daily
std.

\begin{table}[h!]
\centering
\caption{Portfolio-level backtest metrics, full sample.}
\begin{tabular}{lcrrrrr}
\toprule
\textbf{Strategy} & \textbf{Wt.} & \textbf{Cum.\ Ret.} & \textbf{Ann.\ Vol.} & \textbf{Sharpe} & \textbf{Max DD} & \textbf{Turnover} \\
\midrule
Top-10 momentum (lb=30)         & Unif & 2.372 & 0.276          & \textbf{1.022} & $-0.430$           & 0.328 \\
Risk-adj top-10 momentum        & IV   & 1.187 & \textbf{0.181} & 0.961          & $\mathbf{-0.275}$  & 0.443 \\
SMA crossover (20/50)           & Unif & 1.964 & 0.357          & 0.752          & $-0.255$           & \textbf{0.082} \\
Equal-weight buy-and-hold       & Unif & 0.930 & 0.207          & 0.741          & $-0.310$           & 0.014 \\
\bottomrule
\end{tabular}
\end{table}

\begin{figure}[h!]
\centering
\includegraphics[width=0.78\linewidth]{benchmark_nav.png}
\caption{Portfolio NAV. Uniform top-10 momentum dominates on cumulative growth.}
\label{fig:port-nav}
\end{figure}

\begin{figure}[h!]
\centering
\includegraphics[width=0.78\linewidth]{benchmark_drawdown.png}
\caption{Portfolio drawdowns. Inverse-volatility weighting nearly halves the worst drawdown of top-10 momentum.}
\label{fig:port-dd}
\end{figure}

\textbf{Interpretation.} On the same selection rule, uniform weighting
edges inverse-vol on Sharpe (1.02 vs.\ 0.96) and decisively wins on
cumulative return ($2.37\times$ vs.\ $1.19\times$); inverse-vol cuts
volatility by about a third (27.6\% $\to$ 18.1\%) and the drawdown by
roughly half ($-43\%$ $\to$ $-27\%$). The reason uniform wins on
Sharpe in this sample is regime-specific: the highest-momentum names in
2021--2026 are also the highest-volatility mega-cap tech names, so
inverse-vol weighting under-allocates to exactly the winners. This is
the empirical motivation for the Deliverable-5 strategies that
\emph{concentrate} into the strongest momentum names rather than dilute
them via vol scaling.

% ===============================================================
\section{Deliverable 5 -- Benchmarks vs.\ Two New Strategies}

\textbf{Benchmark 1 (SMA crossover).} Equal-weight Nasdaq-100 stocks
whose 20-day SMA exceeds the 50-day SMA; cash if none qualify.
\texttt{SMACrossoverPortfolioStrategy(20, 50)}.

\textbf{Benchmark 2 (Top-K momentum).} Equal-weight the top 10 names
by trailing 30-day return. \texttt{TopKMomentumStrategy(30, 10)}.

\textbf{New strategy 1 -- ConcentratedMomentum.} Top-$K$ by raw trailing
$\ell$-day return, restricted to names with strictly positive trailing
return, weighted equal- or inverse-vol (chosen by the grid). Concentrating
into a small $K$ raises the loading on the strongest momentum signal; the
positive-return gate parks the portfolio in cash during broad-market
downturns. \texttt{ConcentratedMomentumStrategy(lookback, k, weighting,
volatility\_window)}.

\textbf{New strategy 2 -- MomentumWithDrawdownControl.} Same top-$K$
positive-return selection, equal-weighted; if the equal-weight market
index has drawn down by more than $|\text{thr}|$ over the last
\texttt{market\_window} days, gross exposure is scaled to
\texttt{defensive\_exposure} (residual in cash). The drawdown signal
uses only past returns.
\texttt{MomentumWithDrawdownControlStrategy(lookback, k,
drawdown\_threshold, defensive\_exposure, market\_window)}.

\textbf{Protocol.} Tune each new strategy on the first 60\% of dates
(2021-04-13 $\to$ 2024-04-09) by maximizing in-sample Sharpe over a
small grid; evaluate the tuned strategies and the two benchmarks on the
held-out 40\% test window (2024-04-10 $\to$ 2026-04-10, 502 days).
Grid for ConcentratedMomentum:
\texttt{lookback}$\in\{20,30,45,60,90\}$, $k\in\{3,5,7,10,15\}$,
\texttt{weighting}$\in$\{equal, inverse\_vol\}, \texttt{vol\_window}=20.
Grid for MomentumWithDrawdownControl:
\texttt{lookback}$\in\{20,30,45,60,90\}$, $k\in\{3,5,7,10\}$,
\texttt{drawdown\_threshold}$\in\{-0.05,-0.10,-0.15\}$,
\texttt{defensive\_exposure}=0.5, \texttt{market\_window}=120. Both grid
winners selected $k=3$ and $\ell=90$ (Concentrated chose
\texttt{inverse\_vol}; DD-control chose \texttt{thr}$=-0.10$).

\begin{table}[h!]
\centering
\caption{Benchmarks vs.\ new strategies, held-out test window (2024-04-10 $\to$ 2026-04-10, 502 days).}
\begin{tabular}{lrrrrrr}
\toprule
\textbf{Strategy} & \textbf{Cum.\ Ret.} & \textbf{Ann.\ Ret.} & \textbf{Ann.\ Vol.} & \textbf{Sharpe} & \textbf{Max DD} & \textbf{Turnover} \\
\midrule
\textbf{ConcentratedMom} ($k{=}3$, $\ell{=}90$, IV)         & 4.103 & 1.266 & 0.477 & \textbf{1.951} & $-0.397$ & 0.236 \\
\textbf{MomDDCtrl} ($k{=}3$, $\ell{=}90$, thr=$-10\%$)      & 3.843 & 1.207 & 0.465 & \textbf{1.937} & $-0.392$ & 0.173 \\
B2: Top-10 momentum (lb=30)                                  & 1.246 & 0.501 & 0.279 & 1.594          & $-0.288$ & 0.302 \\
B1: SMA crossover (20/50)                                    & 0.306 & 0.143 & 0.167 & 0.884          & $\mathbf{-0.181}$ & \textbf{0.069} \\
\bottomrule
\end{tabular}
\end{table}

\begin{figure}[h!]
\centering
\includegraphics[width=0.78\linewidth]{benchmarks_vs_new_nav.png}
\caption{Test-window NAV. Both new strategies finish well above the top-K benchmark.}
\label{fig:new-nav}
\end{figure}

\begin{figure}[h!]
\centering
\includegraphics[width=0.78\linewidth]{benchmarks_vs_new_drawdown.png}
\caption{Test-window drawdowns. The new strategies pay for upside with deeper drawdowns ($\approx -40\%$) than top-K ($-29\%$).}
\label{fig:new-dd}
\end{figure}

\textbf{Result.} Both new strategies beat both benchmarks on Sharpe.
ConcentratedMomentum at 1.951 and MomDDCtrl at 1.937 clear the top-K
benchmark (1.594) by 0.34--0.36 and the SMA benchmark (0.884) by more
than a full unit. The verdict is recorded in
\texttt{outputs/logs/new\_strategies\_log.json::comparison\_vs\_benchmarks}
with \texttt{beats\_both\_benchmarks: true} for both strategies.

\textbf{Caveats.} Three caveats temper the read. (i) The Sharpe gain
comes mostly from amplified \emph{return} ($\approx 2.5\times$ top-K's
ann.\ ret.) and only partly from a worse \emph{volatility} trade
($\approx 1.7\times$); MDD is correspondingly deeper ($-40\%$ vs.
$-29\%$). Higher Sharpe is not lower risk. (ii) Both grid winners sit at
the boundary of the search ($k=3$ smallest, $\ell=90$ longest), a
structural overfitting flag --- the true optimum may lie outside the
grid. (iii) On the test window the standard error of an annualized
Sharpe of 1.95 with $N=502$ is
$\sqrt{(1+\tfrac{1}{2}\widehat{\mathrm{SR}}^2)/N}\cdot\sqrt{T}\approx 0.71$,
so the 95\% CI on the new strategies overlaps with top-K's CI; beating
the top-K benchmark by 0.34 Sharpe is real evidence but not statistically
overwhelming. The strategies should be re-validated under walk-forward
retuning, realistic transaction costs, and market regimes other than the
2024--2026 AI rally before being treated as production-ready.

% ===============================================================
\section{Extensions: UCB Meta-Strategy and Risk Diagnostics}
\label{sec:ext}

\textbf{UCB-1 over four arms.} \texttt{UCBMetaStrategy(arms, c=1.0)}
selects each day by
$\mathrm{UCB}_k(t) = \overline{r}_k + c\sqrt{\ln T(t)/n_k(t)}$
across SMA crossover, top-10 momentum, risk-adjusted momentum
($\ell{=}60$), and low-volatility momentum ($\ell{=}60$); the chosen
arm's weights are forwarded and \texttt{update\_after\_return} credits
that arm's reward after the realized $R_{p,t+1}$. Selection counts on
the full sample are nearly uniform (315/315/313/311 of 1{,}254 trials)
because the per-day reward differences are sub-basis-point and the
exploration bonus dominates.

\begin{table}[h!]
\centering
\caption{UCB meta-strategy and arms, full sample.}
\begin{tabular}{lrrrrr}
\toprule
\textbf{Strategy} & \textbf{Cum.\ Ret.} & \textbf{Ann.\ Vol.} & \textbf{Sharpe} & \textbf{Max DD} & \textbf{Calmar} \\
\midrule
Top-10 momentum (lb=30)                & 2.372 & 0.276 & \textbf{1.022} & $-0.430$           & 0.642          \\
Risk-adj momentum ($\ell{=}60$)        & 0.998 & 0.169 & 0.908          & $\mathbf{-0.200}$  & 0.747          \\
\textbf{UCB meta} ($c{=}1$, $n{=}4$)   & 2.079 & 0.361 & 0.768          & $-0.234$           & \textbf{1.084} \\
SMA crossover (20/50)                  & 1.964 & 0.357 & 0.752          & $-0.255$           & 0.957          \\
Low-vol momentum ($\ell{=}60$)         & 0.399 & 0.149 & 0.528          & $-0.215$           & 0.325          \\
\bottomrule
\end{tabular}
\end{table}

UCB does not improve Sharpe over the best individual arm but cuts max
drawdown by $\sim 20$pp vs.\ top-K and posts the best Calmar in the
table. The cost is $\sim 4\times$ higher turnover; under any realistic
transaction cost the Sharpe gap would widen against UCB.

\textbf{Tail risk and Sharpe uncertainty.} \texttt{src/risk.py} computes
empirical 5\%-VaR / 5\%-CVaR, an i.i.d.\ bootstrap of NAV paths, and
an approximate normal CI for the annualized Sharpe under i.i.d.-Gaussian
returns,
$\mathrm{SE}(\widehat{\mathrm{SR}}) \approx
\sqrt{(1 + \tfrac{1}{2}\widehat{\mathrm{SR}}^2)/N}\cdot\sqrt{T}$.
This is not a full Lo (2002) HAC adjustment; the intervals widen if
returns are autocorrelated or fat-tailed. Top-10 momentum has the worst
tail (CVaR $-3.82\%$, MDD $-43\%$); risk-adjusted momentum
($\ell{=}60$) has the cleanest tail (CVaR $-2.51\%$, MDD $-20\%$). On
the full sample only top-10 momentum's 95\% Sharpe CI ($[0.14, 1.90]$)
sits strictly above zero; the other arms' CIs all include zero, so
on this sample length differences smaller than $\sim 0.5$ Sharpe are
statistically indistinguishable from noise.

\begin{figure}[h!]
\centering
\includegraphics[width=0.72\linewidth]{bootstrap_nav_paths.png}
\caption{Bootstrap NAV paths for top-10 momentum. With identical mean and variance, the 5th--95th percentile spread of terminal NAV is roughly $8\times$ ($1.18 \to 9.18$); much of any realized track record is sequence luck.}
\label{fig:bootstrap}
\end{figure}

% ===============================================================
\section{Limitations and Reproducibility}

\textbf{Limitations.} Daily-close fillability with no spread or
slippage; zero default transaction cost (high-turnover strategies ---
UCB at 1.40, the new $k{=}3$ momentum at 0.17--0.24, top-K at 0.30 ---
would lose materially more alpha than the SMA portfolio at 0.07 under
realistic costs); survivorship bias in the constituent panel;
single 60/40 split with boundary-parameter selection in Deliverable 5; a
single dominant macro regime (post-COVID rates $\to$ AI mega-cap rally)
in the 5-year window; Sharpe CIs assume i.i.d.\ returns.

\textbf{Reproducibility.} Place the dataset at
\texttt{data/nasdaq100\_daily\_5y.csv} and run from the project root:

\begin{small}
\begin{verbatim}
python -m venv .venv && source .venv/bin/activate
pip install -r requirements.txt

python experiments/run_all.py            # engine smoke test
python experiments/run_single_stock.py   # D3
python experiments/run_benchmarks.py     # D4
python experiments/run_new_strategies.py # D5
python experiments/run_ucb_strategy.py   # extension
python experiments/run_risk_analysis.py  # extension
\end{verbatim}
\end{small}

(On macOS, set \texttt{MPLBACKEND=Agg} before running plot-producing
scripts to avoid backend autodetection.) Each script writes a CSV
metrics table, NAV/drawdown PNGs, and a JSON log to
\texttt{outputs/tables/}, \texttt{outputs/figures/},
\texttt{outputs/logs/}. The Deliverable-5 verdict is computed in
\texttt{outputs/logs/new\_strategies\_log.json}.

\end{document}
